\documentclass[11pt]{article}

\usepackage[margin=1in]{geometry}
\usepackage{booktabs}
\usepackage{hyperref}
\usepackage{amsmath}
\usepackage{amssymb}
\usepackage{xcolor}

\hypersetup{
  colorlinks=true,
  linkcolor=blue!45!black,
  urlcolor=blue!45!black,
  citecolor=blue!45!black
}

\newcommand{\TzEL}{TzEL}
\newcommand{\tez}{\ensuremath{\mathrm{\text{ꜩ}}}}
\newcommand{\mutez}{\ensuremath{\mu\tez}}
\newcommand{\Hash}{\mathsf{H}}
\newcommand{\cm}{\mathsf{cm}}
\newcommand{\nf}{\mathsf{nf}}
\newcommand{\rcm}{\mathsf{rcm}}

\title{\TzEL: A Lite Whitepaper\\
\large Post-Quantum Private Transactions on Tezos Smart Rollups}
\author{TzEL contributors}
\date{\today}

\begin{document}
\maketitle

\begin{abstract}
\TzEL{} is a shielded-payment protocol for Tezos smart rollups. It keeps the
standard private-UTXO architecture of notes, commitments, nullifiers, and
historical Merkle roots, but instantiates the privacy path without primitives
that large-scale quantum computers are expected to break. Notes are encrypted
with ML-KEM-based memo keys; spends are authorized with one-time hash-based
signatures verified inside the zero-knowledge proof; and proofs are recursive
STARKs. The design target is a shielded-payment rollup whose archived note
ciphertexts and spend authorizations do not depend on discrete-log hardness.
\end{abstract}

\section{Scope}

This is a compact design paper. The normative protocol specification, exact
encodings, constants, hash personalizations, and test vectors are in
\texttt{specs/spec.md}. Security notes and caveats are in
\texttt{specs/security.md}.

\section{Motivation}

Privacy-preserving payment systems on blockchains publish ciphertexts as part
of transaction data. Those ciphertexts can be archived indefinitely. If the
encryption protecting them is later broken, old note payloads and memos can be
decrypted retroactively. This is the private-payment form of
harvest-now-decrypt-later.

Hiding values today is therefore not enough when the ciphertexts are protected
by assumptions that large-scale quantum computers are expected to break.
\TzEL{} asks whether the shielded-payment model can be instantiated end to end
with components chosen for that threat model, while still fitting the execution
and data costs of a Tezos rollup.

\section{Protocol at a Glance}

\TzEL{} is an append-only private UTXO system. The rollup maintains:

\begin{itemize}
  \item a commitment tree of note commitments;
  \item a set of spent nullifiers;
  \item a set of historical Merkle roots used as spend anchors;
  \item public bridge balances and bridge configuration;
  \item verifier configuration and fee-policy state.
\end{itemize}

There are three private transaction families:

\begin{center}
\begin{tabular}{lll}
\toprule
Transaction & Inputs & Outputs \\
\midrule
Shield & public balance & user note, DAL-producer note \\
Transfer & 1--7 private notes & recipient note, change note, DAL-producer note \\
Unshield & 1--7 private notes & public credit, optional change note, DAL-producer note \\
\bottomrule
\end{tabular}
\end{center}

Every shield, transfer, and unshield pays two separate resource prices. The
first is a burned rollup fee, used to price rollup execution congestion. The
second is a private DAL-producer fee note, used to pay the party publishing the
large payload into DAL.

\section{Cryptographic Profile}

\begin{center}
\begin{tabular}{lll}
\toprule
Role & Instantiation & Purpose \\
\midrule
Commitments & BLAKE2s-derived field hashes & Bind value, address, and owner tag \\
Nullifiers & BLAKE2s-derived field hashes & Prevent double spends \\
Memo encryption & ML-KEM-768 + ChaCha20-Poly1305 & Encrypt note payloads and memos \\
Detection & ML-KEM-768 clue + short tag & Watch-only candidate filtering \\
Spend authorization & WOTS-style hash signatures & One-time authorization per spent note \\
Zero knowledge & Cairo AIR + recursive STARK & Prove spend validity and authorization \\
\bottomrule
\end{tabular}
\end{center}

The protocol does not claim that commitments, nullifiers, or Merkle trees are
new. They are the standard shielded-UTXO skeleton. The important substitution is
that the encryption, authorization, and proof systems are chosen to avoid
assumptions, such as discrete logarithms, that large-scale quantum computers
are expected to break.

\paragraph{Why these primitives.}
The cryptographic choices follow from the public permanence of shielded-payment
data. Memo encryption must remain private even if ciphertexts are archived for
many years, hence ML-KEM is used for the key encapsulation layer. Spend
authorization must not rely on reusable signatures that are vulnerable to
quantum discrete-log attacks, hence one-time hash signatures are used and
consumed through an authorization tree. The proof system must verify private
state transitions without introducing a pairing-based verification assumption,
hence recursive STARKs are used for the proof path. Hash functions then provide
the common substrate for commitments, nullifiers, Merkle nodes, address binding,
and one-time-signature chains.

\section{Keys, Capabilities, and Addresses}

The wallet derives independent spending and incoming branches from a master
secret. The spending branch contains nullifier and authorization material. The
incoming branch contains diversifier, viewing, and detection material.

\[
\begin{aligned}
spend\_seed &= \Hash(\mathrm{TAG\_SPEND}, master\_sk) \\
incoming\_seed &= \Hash(\mathrm{TAG\_INCOMING}, master\_sk)
\end{aligned}
\]

This split is deliberate. A user may want to delegate scanning without
delegating spending. They may also want to detect likely incoming notes without
letting the scanning service read memo contents. Separating the branches makes
those reduced-capability exports possible.

For address index \(j\), the wallet derives:

\begin{itemize}
  \item \(d_j\), the address diversifier;
  \item \(nk\_spend_j\), the per-address nullifier secret;
  \item \(nk\_tag_j = \Hash_{\mathrm{nktg}}(nk\_spend_j)\), the public
        nullifier binding tag;
  \item \(auth\_root_j, pub\_seed_j\), the public authorization-tree key;
  \item ML-KEM viewing and detection keypairs.
\end{itemize}

The address diversifier \(d_j\) lets one wallet derive many unlinkable receiving
addresses. The per-address nullifier secret \(nk\_spend_j\) prevents a single
global nullifier key from being exposed to a prover. The public \(nk\_tag_j\)
lets the sender bind a note to the correct nullifier derivation without learning
the nullifier secret. The authorization root and public seed identify the
address's one-time spend-authority tree. The two ML-KEM keys separate note
viewing from cheap note detection.

Viewing and detection keys are derived deterministically from the incoming
branch and the address index. This makes wallet recovery and watch-only export
well defined: the same master material and index reproduce the same ML-KEM
keypairs, while different address indices produce independent receiving
addresses.

A payment address is:

\[
address_j =
(d_j, auth\_root_j, pub\_seed_j, nk\_tag_j, ek^v_j, ek^d_j)
\]

The capability hierarchy is:

\begin{center}
\begin{tabular}{ll}
\toprule
Capability & What it can do \\
\midrule
Detection & Flag likely incoming notes with false positives \\
Incoming view & Decrypt incoming memos and recover received notes \\
Full view & Track spent/unspent status for known address metadata \\
Spend & Authorize transactions with nullifier and auth material \\
\bottomrule
\end{tabular}
\end{center}

Incoming viewing alone does not reconstruct the authorization root or nullifier
binding metadata for arbitrary historical addresses. A wallet accepts a note
only after matching local address metadata and recomputing the note commitment.

\paragraph{Detection keys.}
The detection key exists for scalability. A wallet or detection service should
not have to try full memo decryption against every published note, and a
watch-only service should not need the memo viewing key. The sender therefore
creates a detection ciphertext and a short tag under \(ek^d_j\). A detector
holding the corresponding detection secret can cheaply test whether a note is a
candidate. The result is intentionally fuzzy: true notes are detected, unrelated
notes may false-positive, and malicious senders can omit or corrupt the clue.
Detection narrows the scan set; it is not note acceptance.

\paragraph{Viewing keys.}
The viewing key decrypts the note payload \(v \parallel rseed \parallel memo\).
It is separated from detection because reading memos is a strictly stronger
capability than flagging candidate notes. After decryption, the wallet still
recomputes the commitment using local address metadata; decryption alone is not
enough to accept funds.

\section{Notes and Owner Binding}

A note commitment binds the recipient address, value, note randomness, and a
fused owner tag. For note value \(v\), note seed \(rseed\), and address
metadata \((d_j, auth\_root_j, pub\_seed_j, nk\_tag_j)\):

\[
  \rcm = \Hash(\Hash(\mathrm{TAG\_RCM}), rseed)
\]

\[
  owner\_tag_j =
  \Hash_{\mathrm{owner}}(auth\_root_j, pub\_seed_j, nk\_tag_j)
\]

\[
  \cm =
  \Hash_{\mathrm{commit}}(d_j, v, \rcm, owner\_tag_j)
\]

The owner tag is important. Without it, an attacker could try to spend a
commitment under unrelated nullifier material. The chain

\[
nk\_spend_j \rightarrow nk\_tag_j \rightarrow owner\_tag_j \rightarrow \cm
\]

forces the commitment to be consistent with the nullifier key material that is
later used in the spend proof.

This also explains why the payment address includes both authorization material
and nullifier-binding material. A sender needs enough public information to
construct a note that the recipient can later spend, but not enough information
to compute the recipient's nullifiers or spend authorizations.

\section{Position-Dependent Nullifiers}

The nullifier includes both the commitment and the leaf position:

\[
  \nf =
  \Hash_{\mathrm{nf}}(nk\_spend_j,
    \Hash_{\mathrm{nf}}(\cm, pos))
\]

This prevents two identical commitments inserted at different leaves from
collapsing to one nullifier. The circuit proves the nullifier computation and
Merkle membership of the commitment at \(pos\). The rollup then checks that the
nullifier has not appeared before and inserts it into the global nullifier set.

\section{One-Time Spend Authorization}

Each address owns an XMSS-style tree of one-time WOTS public keys. The default
authorization depth is 16, giving \(2^{16}\) one-time keys per address. A spend
uses exactly one unused key index.

The WOTS parameters are:

\begin{itemize}
  \item base \(w = 4\);
  \item 133 chains: 128 message digits and 5 checksum digits;
  \item chain length \(w - 1 = 3\).
\end{itemize}

The wallet signs a transaction sighash with the one-time secret key. Inside the
STARK, the circuit:

\begin{enumerate}
  \item recomputes the transaction sighash from public outputs;
  \item decomposes it into base-4 WOTS digits;
  \item hashes each signature chain forward to recover the WOTS public key;
  \item compresses the recovered public key with the address L-tree;
  \item proves membership of that leaf under \(auth\_root_j\).
\end{enumerate}

No authorization leaf, public key, or signature is published. The STARK proof is
the public evidence that the spend was authorized.

\paragraph{Why hide the authorization leaf.}
Publishing the one-time public key or auth-tree leaf would create an additional
public handle for linking spends to address material. Instead, the proof shows
that a valid one-time key exists under the private note owner's authorization
root, while the public outputs reveal only the nullifiers and commitments needed
by the rollup.

\paragraph{Why one-time signatures.}
WOTS-style signatures are hash-based and simple to verify in a STARK, but they
are one-time. Reusing a key leaks chain preimages that can enable forgery. This
is why the wallet must advance and persist the next authorization index before
submission, and why address rotation is required before the finite auth tree is
exhausted.

\section{Delegated Proving}

The prover is not trusted to choose transaction effects. The wallet constructs
the outputs, computes the sighash, signs it with the relevant one-time keys,
and gives the prover only the witness needed to prove the statement. Because
the WOTS verification happens inside the circuit over the public outputs, a
prover cannot change recipients, nullifiers, commitments, fees, withdrawal
targets, or memo hashes without invalidating the proof.

The public statement intentionally contains only what the rollup needs:
historical root, nullifiers, fees, output commitments, public account binding
fields, and memo hashes. Spend authorization material remains private witness.

\section{Sighash and Replay Binding}

For transfer and unshield, the one-time signature is over a circuit-specific
sighash derived from the public outputs. The circuit recomputes this sighash
internally before checking the WOTS signature. The sighash includes:

\begin{itemize}
  \item a transaction-type tag;
  \item the deployment authorization domain;
  \item the input root;
  \item all nullifiers;
  \item public values such as fee, public withdrawal amount, and recipient id;
  \item output commitments and memo hashes.
\end{itemize}

The transaction-type tag prevents a signature for one circuit from being reused
in another circuit. The authorization domain separates deployments, forks, and
verifier migrations that might otherwise share a root history. Including
commitments, public account identifiers, fees, and memo hashes prevents an
untrusted prover or relayer from changing the transaction effects after the
wallet authorizes them.

There is no separate public signature verification step for private spends.
The proof itself is the authorization check.

\section{Transaction Statements}

\subsection{Shield}

Shield moves public balance into two private notes: the user's note and the
DAL-producer fee note. The public outputs are:

\[
[v_{pub}, fee, producer\_fee, \cm_{user}, \cm_{producer},
sender\_id, mh_{user}, mh_{producer}]
\]

The circuit proves both commitments and enforces \(producer\_fee > 0\). The
rollup checks sender binding, memo hashes, proof validity, and the current
required burned fee. The sender is debited by:

\[
v_{pub} + fee + producer\_fee
\]

Shield has no private input note, so it does not need a Merkle path, nullifier,
or WOTS spend authorization. Instead, it must bind the proved public
\(sender\_id\) to the submitted public account. In the reference environment,
the account string is hashed into a felt. A production verifier environment
would replace that demo binding with a chain-native account serialization and
check.

\subsection{Transfer}

Transfer consumes \(N\) private notes, where \(1 \leq N \leq 7\), and creates
recipient, change, and producer-fee notes. The public outputs are:

\[
[auth\_domain, root, \nf_0,\ldots,\nf_{N-1}, fee,
\cm_1, \cm_2, \cm_3, mh_1, mh_2, mh_3]
\]

The circuit proves input note commitments, Merkle paths, nullifiers, WOTS
authorization, output commitments, pairwise nullifier distinctness, and:

\[
\sum_{i=0}^{N-1} v_i = v_1 + v_2 + v_3 + fee
\]

where \(v_3\) is the DAL-producer fee note and must be positive.

All input nullifiers are constrained pairwise distinct inside the circuit. The
rollup additionally checks them against the global nullifier set. All values are
range-checked as 64-bit quantities and arithmetic is performed with enough
width to avoid wraparound.

The input count \(N\) is not private. It is visible from the number of
nullifiers in the public output list.

\subsection{Unshield}

Unshield consumes \(N\) private notes, releases public value, and creates a
producer-fee note plus optional private change. The public outputs are:

\[
[auth\_domain, root, \nf_0,\ldots,\nf_{N-1}, v_{pub}, fee,
recipient\_id, \cm_{change}, mh_{change}, \cm_{fee}, mh_{fee}]
\]

The circuit proves the same spend authorization and input validity as transfer,
then enforces:

\[
\sum_{i=0}^{N-1} v_i =
v_{pub} + v_{change} + v_{producer} + fee
\]

If there is no change output, the change commitment, change memo hash, and
change witness data are constrained to zero.

Constraining the unused change witness prevents prover malleability: a prover
cannot hide nonzero private change data behind a zero public commitment.

\section{Fees}

The burned fee is a rollup-level congestion price. It is deducted from consumed
value and recreated nowhere. The current fee policy has a floor of
\(100000\) \mutez{}; the first two accepted private transactions at an inbox
level pay the floor, and each additional accepted private transaction at that
same level doubles the required fee up to a cap. When the inbox level advances,
the counter resets.

The DAL-producer fee is separate. It is represented as an ordinary shielded
output note. A DAL publisher can choose whether to include a payload based on
whether it can detect and decrypt a producer-fee note paying its configured
address. This is not a consensus rule; it is an operator admission policy.

\paragraph{Why the two fees are separate.}
The burned fee and producer fee price different resources. The burned fee
protects rollup nodes from an execution backlog; no party receives it. The
producer fee pays for scarce DAL publication capacity; it must be recoverable by
the publisher and therefore appears as an ordinary private note output. Keeping
the producer fee as a normal output avoids adding special recipient semantics to
consensus.

\section{Encrypted Notes, Memos, and Detection}

Each output note carries:

\begin{itemize}
  \item the commitment \(\cm\);
  \item an ML-KEM detection ciphertext \(ct_d\);
  \item a short detection tag;
  \item an ML-KEM viewing ciphertext \(ct_v\);
  \item a derived AEAD nonce;
  \item encrypted note plaintext \(v \parallel rseed \parallel memo\).
\end{itemize}

The encrypted plaintext includes a fixed-size user memo. The memo can carry
payment references, return addresses, human-readable text, or application data.
It is padded so the ciphertext does not reveal the memo length.

Detection lets a watch-only service test whether a note is probably relevant
without decrypting the memo. It is only a candidate filter: false positives are
expected, and malicious senders can make detection fail. The wallet accepts a
note only by decrypting it, recomputing the commitment from local address
metadata, and checking equality with the posted commitment.

The circuit does not decrypt memos. Instead, the wallet hashes the posted note
data and includes the resulting memo hash in the public outputs. The rollup
recomputes the hash over posted note data, preventing a relayer from swapping or
stripping ciphertext fields after the proof is produced.

\paragraph{Why ciphertext correctness is outside the circuit.}
The proof binds commitments and memo hashes, but it does not prove that the
viewing ciphertext decrypts to the same \((v, rseed, memo)\) used in the
commitment. That check belongs to the recipient wallet. This keeps the circuit
focused on consensus-critical spend validity while still giving wallets a clear
acceptance rule: decrypt, recompute the commitment from local address metadata,
and reject if it does not match.

\section{Consensus Checks Outside the STARK}

The STARK proves private consistency. The rollup still enforces public
consensus rules:

\begin{itemize}
  \item the proof verifies under the expected program hash;
  \item \(auth\_domain\) matches the configured deployment domain;
  \item the input root is a known historical root;
  \item nullifiers are globally fresh;
  \item output commitments match posted note data;
  \item memo hashes match posted note ciphertexts;
  \item shield sender and unshield recipient identifiers match submitted
        public account data;
  \item the burned fee meets the current rollup fee policy;
  \item appending all output commitments succeeds atomically.
\end{itemize}

These checks are part of the protocol. For example, accepting a proof without
checking the historical root would allow spends against fake local trees, and
accepting memo data without checking the memo hash would allow relayers to
mutate note ciphertexts.

\paragraph{Root validation.}
The root in a spend proof must be one of the rollup's historical roots. Without
this check, an attacker could prove a spend against a private tree that contains
self-created notes.

\paragraph{Executable binding.}
The verifier environment must know which circuit was proved. A proof that some
Cairo task ran is insufficient; the rollup checks that the proved executable
hash matches the configured shield, transfer, or unshield program hash.

\paragraph{Public account identifiers.}
Shield and unshield cross the private/public boundary, so the proof includes a
public account identifier. The reference implementation hashes UTF-8 account
strings into field elements. That is a testing and deployment convention, not a
universal account model; any deployment using chain-native addresses must define
the exact byte encoding.

\paragraph{Tree append ordering and atomicity.}
Output commitments are appended in order, and the new root is recorded after the
transaction. Multi-output transactions must preflight capacity so a failed append
cannot leave the tree, nullifier set, or public balances partially updated.

\section{Rollup and DAL Path}

Recursive proof payloads are roughly 300 KB, so they are routed through DAL.
The operator publishes large payload chunks to DAL and injects a small rollup
inbox pointer on L1. The kernel resolves the pointer, decodes the payload, and
applies the same state transition as direct messages.

Rollup configuration messages can also exceed the direct L1 inbox limit. The
configuration path therefore uses the same DAL pointer mechanism, but with an
extra operator-side validation step: verifier and bridge configuration payloads
are authenticated before publication because they are privileged admin actions.

\section{Wire Formats and Interoperability}

The protocol uses canonical binary encodings based on Tezos Data Encoding for
core objects such as payment addresses, encrypted notes, published notes, and
note-feed items. JSON exposed by the reference command-line tools is a
convenience transport, not the normative wire format.

All field elements are 32-byte little-endian felt252 values. Hash outputs are
BLAKE2s-256 truncated into the felt252 range. Hash roles are separated by
personalization strings or explicit domain-tag constants, as specified in the
normative spec.

\paragraph{Merkle tree structure.}
The commitment and authorization trees are left-right binary Merkle trees.
Leaves are raw field elements, internal nodes are domain-separated BLAKE2s
hashes, and zero nodes are derived recursively. Commitment positions and
authorization key indices are range-checked so that an out-of-range index cannot
alias a valid path.

\paragraph{Domain separation.}
The same hash family is used for many roles: commitments, nullifiers, Merkle
nodes, owner tags, WOTS chains, sighashes, memo hashes, and key derivation.
Domain separation is therefore not cosmetic. It prevents an output from one
role from being reused as a meaningful input in another role.

\section{Transaction Size}

Each encrypted output note is roughly 3.2 KB. The recursive proof dominates
transaction size at roughly 295 KB. Typical payloads are therefore:

\begin{center}
\begin{tabular}{lll}
\toprule
Transaction & Note data & Approximate total \\
\midrule
Shield & 2 output notes & about 302 KB \\
Transfer & 3 output notes & about 305 KB plus nullifiers \\
Unshield & 1--2 output notes & about 299--303 KB plus nullifiers \\
\bottomrule
\end{tabular}
\end{center}

These sizes are why the DAL path is part of the architecture. Direct L1 inbox
messages remain useful for small pointers and selected small messages, but not
for full recursive proof payloads.

\section{Known Leakage and Limitations}

\begin{itemize}
  \item The number of nullifiers reveals the number of inputs.
  \item Transaction type, ordering, and timing are public.
  \item Detection is honest-sender and only a filtering mechanism.
  \item There is no outgoing viewing key in the current protocol.
  \item Delegated provers receive per-address witness material and can learn
        spent-state information for that address.
  \item Reusing a WOTS one-time key compromises spend authority.
  \item Wallet backup and multi-device state management are security-critical
        because they protect one-time key freshness.
\end{itemize}

\section{Performance Snapshot}

The recursive proof payloads in the reference stack are around 300 KB, with
verification around tens of milliseconds in the current benchmark environment.
Proof generation is heavier, measured in tens of seconds for the benchmark
cases, with peak memory in the tens of GB for large-input transactions. These
numbers are implementation measurements, not protocol limits, but they explain
why the Tezos DAL is part of the system design rather than an afterthought.

\section{Conclusion}

\TzEL{} is a post-quantum instantiation of a shielded-payment rollup.
It keeps the private-UTXO model that has proven useful in earlier systems, but
changes the cryptographic substrate: ML-KEM for note encryption, hash-based
one-time signatures for spend authorization, hash-based commitments and
nullifiers, and recursive STARKs for zero knowledge. The result is an
end-to-end Tezos rollup architecture for private payments whose archived note
ciphertexts and spend authorizations do not rely on discrete-log hardness.

\end{document}
